% Source: Fukuda (2026), Lecture Notes on Mathematics for Economics, Ch. 1
% ("Fundamentals of Set Theory", pp. 1-56), especially 1.1-1.7 and 1.10-1.11.
% External references actually cited in this chapter: Corbae et al. (2009), Le
% Gall (2022), Ok (2007), Rudin (1976).

\chapter{Sets, Relations, Logic, and the Real Number System}\label{ch:sets}

\section{Motivation: Why an Economist Starts with Sets}\label{sec:sets-motivation}

An economic model is a statement about a set of alternatives, a rule for
comparing them, and a subset of them that is feasible. A consumer facing prices
$p\in\R^n_{++}$ and wealth $m>0$ chooses from the budget set
$\budget(p,m)=\{x\in\R^n_+ : \ip{p}{x}\leq m\}$; a firm chooses from a production
set; a player in a game chooses from a strategy set. The comparison rule is a
binary relation on the set of alternatives, and the object economists call a
utility function is a numerical representation of that relation. Every
subsequent chapter presupposes that these three objects --- a set, a relation, a
distinguished subset --- have been given a precise meaning.

This chapter therefore does three things. It fixes set-theoretic language
(\autoref{sec:sets-basic}--\autoref{sec:sets-relations}); it records the forms of
argument used throughout the volume (\autoref{sec:sets-logic}); and it
constructs the real number system axiomatically
(\autoref{sec:sets-reals}), because the single property that distinguishes $\R$
from $\Q$ --- the least upper bound property --- is what makes every existence
theorem in later chapters true. The chapter closes with the first substantive
economic application, the representation of a preference relation by a utility
function (\autoref{sec:sets-preferences}).

\section{Sets, Operations, and Mappings}\label{sec:sets-basic}

\subsection{Sets and their specification}

We take the notion of a set as primitive: a \dfn{set} is a collection of objects,
called its \dfnas{element}{elements} or members. We write $a\in A$ when $a$ is an
element of $A$ and $a\notin A$ otherwise. The set with no elements is the
\dfnas{empty set}{empty set}, written $\emptyset$.
\nidx{$\emptyset$}{$\emptyset$, empty set}

A set may be specified by enumeration, as in $\{a_1,\dots,a_n\}$, or by a
defining property relative to an ambient set, as in
$\{x\in X : P(x)\}$. The ambient set matters: the collection of all $x$ with
$P(x)$, taken without an ambient set, is not in general a set, and the
restriction to subsets of a given set is what avoids the classical paradoxes. All
constructions in this volume take place inside an ambient set that is either
named or clear from context.

\begin{definition}[Inclusion and equality]\label{def:inclusion}
	Let $A$ and $B$ be sets. We say $A$ is a \dfn{subset} of $B$, written
	$A\subseteq B$, if $a\in A$ implies $a\in B$. We say $A=B$ if $A\subseteq B$ and
	$B\subseteq A$. If $A\subseteq B$ and $A\neq B$, then $A$ is a \dfn{proper
		subset} of $B$, written $A\subsetneq B$.
\end{definition}

The definition of equality dictates the standard proof technique: to establish
$A=B$ one establishes two inclusions. This is used so often below that it is
worth naming once.

\begin{definition}[Operations on sets]\label{def:set-ops}
	Let $A,B$ be subsets of an ambient set $X$, and let $(A_i)_{i\in I}$ be a family
	of subsets of $X$ indexed by a set $I$.
	\begin{enumerate}[(i)]
		\item The \dfn{complement} of $A$ in $X$ is
		      $\comp{A}\eqdef\{x\in X : x\notin A\}$.
		\item The \dfn{union} and \dfn{intersection} of the family are
		      \[
			      \bigcup_{i\in I}A_i\eqdef\{x\in X : x\in A_i\ \text{for some }i\in I\},
			      \qquad
			      \bigcap_{i\in I}A_i\eqdef\{x\in X : x\in A_i\ \text{for all }i\in I\}.
		      \]
		\item The \dfn{set difference} is $A\setminus B\eqdef A\cap\comp{B}$.
		\item The \dfn{Cartesian product} of sets $A_1,\dots,A_n$ is
		      $\prod_{i=1}^n A_i\eqdef\{(a_1,\dots,a_n) : a_i\in A_i\ \text{for each }i\}$,
		      the set of ordered $n$-tuples.
		\item The \dfn{power set} of $X$ is $\pow{X}\eqdef\{A : A\subseteq X\}$.
	\end{enumerate}
\end{definition}

The index set $I$ in (ii) is arbitrary: it may be finite, countably infinite, or
uncountable. This generality is not idle. In \autoref{ch:topology} the collection
of open sets is required to be closed under arbitrary unions but only under
finite intersections, and the whole of point-set topology turns on that
asymmetry.

\begin{proposition}[De Morgan's laws]\label{prop:demorgan}
	Let $(A_i)_{i\in I}$ be a family of subsets of $X$ with $I\neq\emptyset$. Then
	\[
		\comp{\left(\bigcup_{i\in I}A_i\right)}=\bigcap_{i\in I}\comp{A_i},
		\qquad
		\comp{\left(\bigcap_{i\in I}A_i\right)}=\bigcup_{i\in I}\comp{A_i}.
	\]
\end{proposition}

\begin{proof}
	For the first identity, $x\in\comp{\left(\bigcup_{i}A_i\right)}$ holds if and
	only if $x\in X$ and it is not the case that $x\in A_i$ for some $i\in I$; that
	is, if and only if $x\in X$ and $x\notin A_i$ for every $i\in I$; that is, if and
	only if $x\in\comp{A_i}$ for every $i$, which is $x\in\bigcap_i \comp{A_i}$.

	The second identity follows from the first applied to the family
	$(\comp{A_i})_{i\in I}$ together with $\comp{(\comp{A})}=A$: indeed
	$\comp{\left(\bigcup_i \comp{A_i}\right)}=\bigcap_i A_i$, and taking complements
	of both sides gives the claim.
\end{proof}

\begin{remark}
	The hypothesis $I\neq\emptyset$ is needed. With $I=\emptyset$ the convention
	$\bigcup_{i\in\emptyset}A_i=\emptyset$ and $\bigcap_{i\in\emptyset}A_i=X$ makes
	both identities hold, but only because the intersection over the empty family is
	defined to be the ambient set. Outside an ambient set that convention has no
	meaning, which is a second reason for insisting on one.
\end{remark}

\subsection{Mappings}

\begin{definition}[Mapping]\label{def:mapping}
	Let $A$ and $B$ be sets. A \dfn{mapping} (or function) $f\colon A\to B$ assigns
	to each $a\in A$ exactly one element $f(a)\in B$. The set $A$ is the
	\dfn{domain} and $B$ the \dfnas{codomain}{codomain}. For $S\subseteq A$ and
	$T\subseteq B$,
	\[
		f(S)\eqdef\{f(a) : a\in S\}\subseteq B,\qquad
		f^{-1}(T)\eqdef\{a\in A : f(a)\in T\}\subseteq A
	\]
	are the \dfn{image} of $S$ and the \dfn{preimage} of $T$. The set $f(A)$ is the
	\dfnas{range}{range of a mapping}.
\end{definition}

The preimage is better behaved than the image, and the difference is the reason
that continuity is defined in \autoref{ch:continuity} through preimages rather
than images.

\begin{proposition}[Images and preimages under set operations]\label{prop:image-preimage}
	Let $f\colon A\to B$, let $(S_i)_{i\in I}$ be subsets of $A$, and let
	$(T_i)_{i\in I}$ be subsets of $B$ with $I\neq\emptyset$. Then
	\begin{enumerate}[(i)]
		\item $f^{-1}\!\left(\bigcup_i T_i\right)=\bigcup_i f^{-1}(T_i)$ and
		      $f^{-1}\!\left(\bigcap_i T_i\right)=\bigcap_i f^{-1}(T_i)$;
		\item $f^{-1}(\comp{T})=\comp{\left(f^{-1}(T)\right)}$ for every $T\subseteq B$;
		\item $f\!\left(\bigcup_i S_i\right)=\bigcup_i f(S_i)$, whereas in general only
		      $f\!\left(\bigcap_i S_i\right)\subseteq\bigcap_i f(S_i)$.
	\end{enumerate}
\end{proposition}

\begin{proof}
	(i) We have $a\in f^{-1}\!\left(\bigcap_i T_i\right)$ if and only if
	$f(a)\in T_i$ for every $i$, if and only if $a\in f^{-1}(T_i)$ for every $i$.
	The union case is identical with ``some'' in place of ``every''.

	(ii) $a\in f^{-1}(\comp{T})$ if and only if $f(a)\notin T$, if and only if
	$a\notin f^{-1}(T)$.

	(iii) If $b\in f\!\left(\bigcup_i S_i\right)$ then $b=f(a)$ for some
	$a\in S_{i_0}$, so $b\in f(S_{i_0})$; conversely each $f(S_i)\subseteq
		f\!\left(\bigcup_j S_j\right)$. For intersections, if $a\in\bigcap_i S_i$ then
	$f(a)\in f(S_i)$ for every $i$, giving the stated inclusion.
\end{proof}

\begin{eg}[Failure of equality for images of intersections]\label{eg:image-intersection}
	Let $f\colon\R\to\R$ be $f(x)=x^2$, $S_1=(-\infty,0)$ and $S_2=(0,\infty)$. Then
	$S_1\cap S_2=\emptyset$, so $f(S_1\cap S_2)=\emptyset$, while
	$f(S_1)=f(S_2)=(0,\infty)$, so $\bigcap_i f(S_i)=(0,\infty)$. The inclusion in
	\autoref{prop:image-preimage}(iii) is therefore strict in general, and it is an
	equality for all families if and only if $f$ is injective.
\end{eg}

\begin{definition}[Injection, surjection, bijection]\label{def:inj-surj}
	A mapping $f\colon A\to B$ is
	\dfnas{injective}{injective mapping} if $f(a)=f(a')$ implies $a=a'$;
	\dfnas{surjective}{surjective mapping} if $f(A)=B$; and
	\dfnas{bijective}{bijective mapping} if it is both. A bijection $f$ has a
	unique inverse $f^{-1}\colon B\to A$ satisfying $f^{-1}\circ f=\id_A$ and
	$f\circ f^{-1}=\id_B$.
\end{definition}

\begin{definition}[Finite, countable, uncountable]\label{def:countable}
	A set $A$ is \dfn{finite} if $A=\emptyset$ or there is a bijection
	$\{1,\dots,n\}\to A$ for some $n\in\N$; \dfnas{countably infinite}{countably
		infinite set} if there is a bijection $\N\to A$; \dfn{countable} if it is finite
	or countably infinite; and \dfn{uncountable} otherwise.
\end{definition}

\begin{proposition}\label{prop:countable-union}
	A countable union of countable sets is countable. The set $\Q$ is countably
	infinite; the set $\R$ is uncountable.
\end{proposition}

\begin{proof}
	The first two assertions are standard diagonal enumerations; see
	\textcite[Thm.~2.12 and its corollary]{rudin1976principles}. For the third,
	Cantor's diagonal argument shows that no map $\N\to(0,1)$ is surjective: given
	$f\colon\N\to(0,1)$ with decimal expansions $f(n)=0.d_{n1}d_{n2}\dots$ (choosing
	the expansion that does not terminate in all nines), the number
	$x=0.c_1c_2\dots$ with $c_n=5$ if $d_{nn}\neq5$ and $c_n=4$ otherwise lies in
	$(0,1)$ and differs from every $f(n)$ in its $n$-th digit. Since
	$(0,1)\subseteq\R$, the set $\R$ is uncountable.
\end{proof}

\begin{remark}[Where countability matters in economics]\label{rem:countability-econ}
	The distinction is not decorative. In \autoref{ch:probability} a probability
	measure can be assigned to every subset of a countable sample space, but on an
	uncountable sample space such as $[0,1]$ no countably additive probability
	measure exists on the full power set that extends length; one must restrict
	attention to a $\sigma$-algebra of measurable sets
	\autocite[Ch.~1]{legall2022measure}. The Gerzensee treatment makes exactly this
	point when it moves from ``an at most countable set $\Omega$''
	\autocite[slide 3]{fukuda2026probability} to the remark that a continuum
	sample space requires a $\sigma$-algebra
	\autocite[slide 37]{fukuda2026analysis}.
\end{remark}

\section{Relations, Orders, and Preferences}\label{sec:sets-relations}

\begin{definition}[Binary relation]\label{def:relation}
	A \dfn{binary relation} $R$ on a set $X$ is a subset $R\subseteq X\times X$; one
	writes $xRy$ for $(x,y)\in R$. The relation is
	\begin{enumerate}[(i)]
		\item \dfnas{reflexive}{reflexive relation} if $xRx$ for all $x\in X$;
		\item \dfnas{complete}{complete relation} if for all $x,y\in X$, $xRy$ or $yRx$
		      (or both);
		\item \dfnas{transitive}{transitive relation} if $xRy$ and $yRz$ imply $xRz$;
		\item \dfnas{symmetric}{symmetric relation} if $xRy$ implies $yRx$;
		\item \dfnas{antisymmetric}{antisymmetric relation} if $xRy$ and $yRx$ imply
		      $x=y$.
	\end{enumerate}
\end{definition}

Completeness implies reflexivity (take $y=x$), so the two are not independent
hypotheses. Economics texts frequently list both; nothing is gained by it, and we
list completeness alone where it is available.

\begin{definition}[Equivalence relation and partition]\label{def:equivalence}
	An \dfn{equivalence relation} on $X$ is a reflexive, symmetric and transitive
	relation, usually written $\sim$. For $x\in X$ the \dfn{equivalence class} of
	$x$ is $[x]\eqdef\{y\in X : y\sim x\}$. A \dfn{partition} of $X$ is a family of
	non-empty, pairwise disjoint subsets whose union is $X$.
\end{definition}

\begin{proposition}\label{prop:equivalence-partition}
	The equivalence classes of an equivalence relation on $X$ form a partition of
	$X$; conversely, every partition of $X$ arises in this way from exactly one
	equivalence relation.
\end{proposition}

\begin{proof}
	Reflexivity gives $x\in[x]$, so the classes are non-empty and cover $X$. Suppose
	$[x]\cap[y]\neq\emptyset$ and take $z$ in the intersection. Then $z\sim x$ and
	$z\sim y$, so by symmetry and transitivity $x\sim y$; hence for any $w\in[x]$ we
	get $w\sim x\sim y$ and $w\in[y]$, and symmetrically, so $[x]=[y]$. Thus distinct
	classes are disjoint.

	Conversely, given a partition $\mathcal{P}$, define $x\sim y$ if $x$ and $y$ lie
	in the same cell. This is reflexive, symmetric and transitive because the cells
	are non-empty, pairwise disjoint, and cover $X$; and its classes are exactly the
	cells. Uniqueness holds because the relation is determined by its classes.
\end{proof}

This proposition is used twice in later chapters: in \autoref{ch:probability}, a
finite partition $(H_j)_{j=1}^m$ of the sample space is precisely the information
structure underlying the law of iterated expectations
\autocite[slide 18]{fukuda2026probability}; and in the theory of choice, the
indifference relation derived from a preference partitions the consumption set
into indifference curves.

\subsection{Preference relations}\label{subsec:preferences}

\begin{definition}[Preference relation]\label{def:preference}
	Let $X$ be a non-empty set of alternatives. A \dfn{preference relation}
	$\succsim$ on $X$ is a complete and transitive binary relation. Its
	\dfn{strict part} $\succ$ and \dfn{indifference part} $\sim$ are defined by
	\[
		x\succ y \iff \bigl(x\succsim y \text{ and not } y\succsim x\bigr),
		\qquad
		x\sim y \iff \bigl(x\succsim y \text{ and } y\succsim x\bigr).
	\]
\end{definition}

\begin{proposition}\label{prop:preference-parts}
	Let $\succsim$ be a preference relation on $X$. Then $\sim$ is an equivalence
	relation, $\succ$ is transitive and irreflexive, and for all $x,y\in X$ exactly
	one of $x\succ y$, $y\succ x$, $x\sim y$ holds.
\end{proposition}

\begin{proof}
	Reflexivity of $\sim$ follows from completeness. Symmetry is immediate from the
	definition. For transitivity of $\sim$, if $x\sim y$ and $y\sim z$ then
	$x\succsim y\succsim z$ gives $x\succsim z$, and $z\succsim y\succsim x$ gives
	$z\succsim x$.

	For $\succ$: irreflexivity is immediate since $x\succsim x$. For transitivity,
	suppose $x\succ y$ and $y\succ z$. Transitivity of $\succsim$ gives
	$x\succsim z$. If also $z\succsim x$, then $z\succsim x\succsim y$ gives
	$z\succsim y$, contradicting $y\succ z$. Hence $x\succ z$.

	For the trichotomy, completeness gives $x\succsim y$ or $y\succsim x$. If both,
	then $x\sim y$; if only the first, $x\succ y$; if only the second, $y\succ x$.
	These three cases are mutually exclusive by construction.
\end{proof}

\begin{definition}[Utility representation]\label{def:utility-rep}
	A function $u\colon X\to\R$ \dfnas{represents}{utility representation} the
	preference relation $\succsim$ if
	\[
		x\succsim y \iff u(x)\geq u(y)\qquad\text{for all }x,y\in X.
	\]
\end{definition}

\begin{theorem}[Representation on a countable set]\label{thm:utility-countable}
	Let $X$ be a countable non-empty set and let $\succsim$ be a preference relation
	on $X$. Then there exists $u\colon X\to\R$ representing $\succsim$, and $u$ may
	be chosen with $u(X)\subseteq(0,1)$.
\end{theorem}

\begin{proof}
	Enumerate $X=\{x_1,x_2,\dots\}$ (finite or infinite) and choose weights
	$w_n>0$ with $W\eqdef\sum_n w_n<\infty$, for instance $w_n=2^{-n}$. Define
	\[
		u(x)\eqdef\frac{1}{W}\sum_{n:\,x\succ x_n}w_n .
	\]
	Suppose $x\succsim y$. If $y\succ x_n$ then $x\succ x_n$: indeed
	$x\succsim y\succsim x_n$ gives $x\succsim x_n$, and were $x_n\succsim x$ we
	would get $x_n\succsim x\succsim y$, contradicting $y\succ x_n$. Hence
	$\{n : y\succ x_n\}\subseteq\{n : x\succ x_n\}$ and $u(x)\geq u(y)$.

	Conversely, suppose $u(x)\geq u(y)$ and, for a contradiction, that
	$y\succ x$. Then as above $\{n : x\succ x_n\}\subseteq\{n : y\succ x_n\}$. The
	inclusion is strict: since $y\succ x$ and $x=x_{n_0}$ for some $n_0$, we have
	$n_0\in\{n : y\succ x_n\}$ while $n_0\notin\{n : x\succ x_n\}$ by irreflexivity
	of $\succ$. As $w_{n_0}>0$, this gives $u(y)>u(x)$, a contradiction. Hence
	$x\succsim y$.

	Finally $0\leq u\leq1$; replacing $u$ by $(1+u)/3$ places the range in $(0,1)$
	without disturbing the representation, since the transformation is strictly
	increasing.
\end{proof}

\begin{remark}[Why countability is not innocuous]\label{rem:lexicographic}
	\autoref{thm:utility-countable} fails on uncountable $X$ without further
	structure. Let $X=\R^2$ with the \dfn{lexicographic preference}
	\[
		(x_1,x_2)\succsim(y_1,y_2)
		\iff
		x_1>y_1,\ \text{or}\ \bigl(x_1=y_1\ \text{and}\ x_2\geq y_2\bigr).
	\]
	This relation is complete and transitive but admits no utility representation.
	Indeed, suppose $u$ represented it. For each $t\in\R$ we have
	$(t,1)\succ(t,0)$, so $u(t,1)>u(t,0)$, and we may pick a rational
	$q(t)\in\bigl(u(t,0),u(t,1)\bigr)$ by \autoref{prop:archimedean}(ii). If
	$t<t'$ then $(t',0)\succ(t,1)$, so $q(t)<u(t,1)<u(t',0)<q(t')$; hence
	$t\mapsto q(t)$ is a strictly increasing injection $\R\to\Q$, contradicting
	\autoref{prop:countable-union}.

	The lesson carries into \autoref{ch:continuity}: a continuous utility
	representation requires the preference to be continuous, in the sense that the
	sets $\{y : y\succsim x\}$ and $\{y : x\succsim y\}$ are closed, and
	lexicographic preference violates exactly this. The general statement is
	Debreu's theorem; see \textcite[Ch.~2]{ok2007real} and
	\textcite[Ch.~4]{corbae2009introduction}.
\end{remark}

\section{Logic and Methods of Proof}\label{sec:sets-logic}

\subsection{Connectives and quantifiers}

For propositions $P,Q$, the implication $P\Rightarrow Q$ is by definition the
proposition $(\neg P)\vee Q$. This is the source of the ``vacuous truth''
convention: $P\Rightarrow Q$ is true whenever $P$ is false, regardless of $Q$.
The truth table is
\[
	\begin{array}{cc|ccc}
		P  & Q  & \neg P & (\neg P)\vee Q & P\Rightarrow Q \\\hline
		\T & \T & \F     & \T             & \T             \\
		\T & \F & \F     & \F             & \F             \\
		\F & \T & \T     & \T             & \T             \\
		\F & \F & \T     & \T             & \T
	\end{array}
\]
which reproduces the table given in \autocite[slide 40]{fukuda2026analysis}. The
convention is used there to show that $\emptyset$ is open: the statement ``if
$x\in\emptyset$ then $\ball{\varepsilon}{x}\subseteq\emptyset$'' is vacuously
true, so $\emptyset$ satisfies the definition of an open set. That argument
recurs in \autoref{ch:topology}.

Quantifier order is not commutative and is a frequent source of error. The
statement
\[
	\forall\varepsilon>0\ \exists\delta>0\ \forall x\in A:\ d(x,x_0)<\delta
	\Rightarrow d(f(x),f(x_0))<\varepsilon
\]
defines continuity of $f$ at $x_0$, where $\delta$ may depend on both
$\varepsilon$ and $x_0$. Moving the quantifier $\forall x_0\in A$ inside, to
\[
	\forall\varepsilon>0\ \exists\delta>0\ \forall x_0\in A\ \forall x\in A:\
	d(x,x_0)<\delta\Rightarrow d(f(x),f(x_0))<\varepsilon,
\]
gives uniform continuity, a strictly stronger property
(\autoref{ch:continuity}). The negation of a quantified statement exchanges
$\forall$ and $\exists$ and negates the matrix:
$\neg\bigl(\forall x\, P(x)\bigr)$ is $\exists x\,\neg P(x)$.

\subsection{Three methods of proof}

\begin{proposition}[Principle of mathematical induction]\label{prop:induction}
	Let $S\subseteq\N$ satisfy $1\in S$ and $n\in S\Rightarrow n+1\in S$. Then
	$S=\N$.
\end{proposition}

The principle is equivalent to the well-ordering of $\N$ --- every non-empty
subset of $\N$ has a least element --- and either may be taken as an axiom.
Strong induction, from ``$\{1,\dots,n\}\subseteq S$ implies $n+1\in S$'' conclude
$S=\N$, follows by applying \autoref{prop:induction} to
$S'=\{n : \{1,\dots,n\}\subseteq S\}$.

\begin{eg}[Induction in a dynamic argument]\label{eg:induction-contraction}
	The error bound in the contraction mapping theorem
	(\autoref{thm:banach-fixed-point}) rests on the estimate
	$d(x_n,x_{n+1})\leq\beta^n d(x_0,x_1)$ for all $n$, which is an induction: the
	case $n=0$ is trivial, and
	$d(x_{n+1},x_{n+2})=d(Tx_n,Tx_{n+1})\leq\beta d(x_n,x_{n+1})\leq
		\beta^{n+1}d(x_0,x_1)$.
\end{eg}

\emph{Proof by contradiction} establishes $P$ by assuming $\neg P$ and deriving a
false statement; \emph{proof by contraposition} establishes $P\Rightarrow Q$ by
establishing $\neg Q\Rightarrow\neg P$, the two being equivalent by the truth
table above. The uniqueness half of the contraction mapping theorem is a proof by
contradiction; \autoref{prop:preference-parts} used contraposition twice.

\begin{remark}[A common error]\label{rem:proof-error}
	``Suppose not; then \dots; hence $P$'' is not a proof by contradiction; it is a
	proof of $P$ with an unused hypothesis. A genuine argument by contradiction must
	terminate in a statement known to be false. Confusing the two conceals which
	hypotheses are actually used, which matters when one later asks whether a
	hypothesis can be weakened --- the recurring question of
	\autoref{ch:optimization}.
\end{remark}

\section{The Real Number System}\label{sec:sets-reals}

The real numbers are characterised, up to isomorphism, by three groups of
axioms. We state them because the third is used in essentially every existence
theorem in this volume.

\begin{definition}[Ordered field]\label{def:ordered-field}
	A \dfn{field} is a set $F$ with operations $+$ and $\cdot$ satisfying the usual
	commutative, associative and distributive laws, with additive identity $0$,
	multiplicative identity $1\neq0$, additive inverses, and multiplicative inverses
	for every element other than $0$. An \dfn{ordered field} is a field with a
	complete transitive antisymmetric relation $\leq$ such that $x\leq y$ implies
	$x+z\leq y+z$ for all $z$, and $0\leq x$, $0\leq y$ imply $0\leq xy$.
\end{definition}

\begin{definition}[Bounds, supremum, infimum]\label{def:sup-inf}
	Let $A\subseteq\R$ be non-empty. An \dfn{upper bound} for $A$ is $b\in\R$ with
	$a\leq b$ for all $a\in A$; $A$ is \dfnas{bounded above}{bounded above} if such
	a $b$ exists. The \dfn{supremum} $\sup A$ is the least upper bound: an upper
	bound $s$ such that $s\leq b$ for every upper bound $b$. \dfn{Lower bound},
	\dfnas{bounded below}{bounded below} and \dfn{infimum} $\inf A$ are defined
	symmetrically. A set bounded above and below is \dfn{bounded}.
	\nidx{$\sup$}{$\sup A$, supremum}\nidx{$\inf$}{$\inf A$, infimum}
\end{definition}

\begin{assumption}[Completeness axiom]\label{ass:completeness-axiom}
	$\R$ is an ordered field in which every non-empty subset that is bounded above
	has a supremum in $\R$.
\end{assumption}

This is the \dfnas{least upper bound property}{least upper bound property}. It
fails in $\Q$: the set $\{q\in\Q : q^2<2\}$ is bounded above in $\Q$ but has no
least upper bound there. The whole difference between $\Q$ and $\R$, and with it
every convergence theorem below, rests on this one axiom.

\begin{proposition}[Characterisation of the supremum]\label{prop:sup-char}
	Let $A\subseteq\R$ be non-empty and bounded above, and let $s\in\R$. Then
	$s=\sup A$ if and only if
	\begin{enumerate}[(i)]
		\item $a\leq s$ for all $a\in A$; and
		\item for every $\varepsilon>0$ there exists $a\in A$ with $a>s-\varepsilon$.
	\end{enumerate}
\end{proposition}

\begin{proof}
	Suppose $s=\sup A$. Then (i) holds since $s$ is an upper bound. For (ii), if
	some $\varepsilon>0$ had $a\leq s-\varepsilon$ for all $a\in A$, then
	$s-\varepsilon$ would be an upper bound smaller than $s$, contradicting
	leastness.

	Conversely, assume (i) and (ii). By (i), $s$ is an upper bound. Let $b$ be any
	upper bound and suppose $b<s$. Applying (ii) with $\varepsilon=s-b>0$ produces
	$a\in A$ with $a>s-(s-b)=b$, contradicting that $b$ is an upper bound. Hence
	$s\leq b$ for every upper bound $b$, so $s$ is least.
\end{proof}

Part (ii) is the form in which the supremum is actually used: it converts the
abstract existence of a least upper bound into the ability to produce elements of
$A$ arbitrarily close to it. Every ``choose $a_n\in A$ with
$a_n>\sup A-1/n$'' argument in later chapters is an application of it.

\begin{proposition}[Archimedean property and density of $\Q$]\label{prop:archimedean}
	\begin{enumerate}[(i)]
		\item For all $x,y\in\R$ with $x>0$ there exists $n\in\N$ with $nx>y$.
		\item For all $x<y$ in $\R$ there exists $q\in\Q$ with $x<q<y$.
	\end{enumerate}
\end{proposition}

\begin{proof}
	(i) Suppose not: then $A=\{nx : n\in\N\}$ is bounded above by $y$. By
	\autoref{ass:completeness-axiom} it has a supremum $s$. Since $x>0$, the number
	$s-x$ is not an upper bound, so $mx>s-x$ for some $m\in\N$, whence
	$(m+1)x>s$, contradicting that $s$ is an upper bound.

	(ii) Since $y-x>0$, part (i) gives $n\in\N$ with $n(y-x)>1$, that is
	$ny>nx+1$. Again by (i) there are integers $m_1>nx$ and $m_2>-nx$, so
	$-m_2<nx<m_1$; let $m$ be the least integer in $\{k\in\Z : -m_2\leq k\leq m_1\}$
	with $k>nx$, which exists by well-ordering. Then $m-1\leq nx$, so
	$nx<m\leq nx+1<ny$, and $q=m/n$ satisfies $x<q<y$.
\end{proof}

\begin{remark}[The extended real line]\label{rem:extended-reals}
	It is convenient to write $\sup A=+\infty$ when $A$ is unbounded above and
	$\inf A=-\infty$ when $A$ is unbounded below, and to set
	$\sup\emptyset=-\infty$, $\inf\emptyset=+\infty$. The resulting object
	$\overline{\R}=\R\cup\{\pm\infty\}$ is not a field --- $\infty-\infty$ is
	undefined --- and no algebraic identity may be applied to it without checking
	that the terms are finite. In optimisation this is not pedantry: the value of an
	infeasible maximisation problem is $-\infty$ by the convention
	$\sup\emptyset=-\infty$, whereas the value of a feasible but unbounded problem
	is $+\infty$, and the two failures call for different remedies.
\end{remark}

\section{Economic Application: Choice, Feasibility, and the Existence Question}
\label{sec:sets-preferences}

We can now state, in the language of this chapter, the problem that occupies the
rest of the volume.

\begin{eg}[The consumer's problem, stated set-theoretically]\label{eg:consumer-setup}
	Let $X=\R^n_+$ be the \dfn{consumption set}, let $\succsim$ be a preference
	relation on $X$, and let prices $p\in\R^n_{++}$ and wealth $m>0$ be given. The
	\dfn{budget set} is
	\[
		\budget(p,m)\eqdef\{x\in\R^n_+ : \ip{p}{x}\leq m\}.
	\]
	The consumer's problem is to find
	\[
		x^\ast\in\budget(p,m)\quad\text{such that}\quad
		x^\ast\succsim x\ \text{for all } x\in\budget(p,m),
	\]
	and the \dfn{demand correspondence} $x^\ast(p,m)$ is the set of all such
	$x^\ast$.
\end{eg}

Four questions arise immediately, and each is answered by a theorem in a later
chapter.

\begin{enumerate}[(a)]
	\item \emph{Does a maximal element exist?} If $\succsim$ is represented by a
	      continuous $u$ and $\budget(p,m)$ is compact, the Weierstrass theorem
	      (\autoref{thm:weierstrass}) gives existence. Compactness of $\budget(p,m)$
	      requires $p\in\R^n_{++}$: if some $p_i=0$ the set is unbounded and existence
	      can fail. This is why the strict positivity of prices is a hypothesis and
	      not a normalisation.
	\item \emph{Is the maximiser unique?} Uniqueness requires strict quasi-concavity
	      of $u$, equivalently strict convexity of preferences
	      (\autoref{ch:convex}).
	\item \emph{Can it be characterised by first-order conditions?} This requires
	      differentiability of $u$ and a constraint qualification
	      (\autoref{ch:optimization}); the Cobb--Douglas case worked in
	      \autocite[slides 21--32]{fukuda2026optimization} is developed there in full.
	\item \emph{How does it move with $(p,m)$?} Continuity of $x^\ast(\cdot,\cdot)$
	      as a correspondence is Berge's maximum theorem
	      (\autoref{thm:berge-maximum}); its differentiability, where it holds,
	      follows from the implicit function theorem
	      (\autoref{thm:implicit-function}).
\end{enumerate}

The point of listing them together is that they are the same four questions ---
existence, uniqueness, characterisation, comparative statics --- asked of every
optimisation problem in economics, including the dynamic ones of
\autoref{ch:dynamic}, where the object being chosen is an infinite sequence and
the ``budget set'' is a set of feasible paths.

\section{Chapter Summary}\label{sec:sets-summary}

Sets are specified relative to an ambient set; equality of sets is proved by
double inclusion; preimages commute with all set operations while images commute
only with unions. A preference relation is a complete, transitive binary
relation; its indifference part partitions the alternatives; and it admits a
utility representation on any countable set, but not in general on an uncountable
one, as lexicographic preference shows. Implication is $(\neg P)\vee Q$, which
makes statements with false hypotheses vacuously true --- the device by which
$\emptyset$ is open. The real numbers form an ordered field with the least upper
bound property; the working form of that property is
\autoref{prop:sup-char}(ii), and the Archimedean property and the density of the
rationals are its first two consequences.

\begin{exercise}\label{ex:sets-1}
	Show that $f\colon A\to B$ is injective if and only if
	$f\left(\bigcap_{i\in I}S_i\right)=\bigcap_{i\in I}f(S_i)$ for every non-empty
	family $(S_i)_{i\in I}$ of subsets of $A$. (One direction is
	\autoref{eg:image-intersection}; for the other, consider two-element families.)
\end{exercise}

\begin{exercise}\label{ex:sets-2}
	Let $\succsim$ be a preference relation on a finite set $X$. Show directly, by
	induction on the number of elements of $X$, that a $\succsim$-maximal element
	exists. Where does the argument break down when $X$ is infinite, and which
	hypothesis of the Weierstrass theorem repairs it?
\end{exercise}

\begin{exercise}\label{ex:sets-3}
	Let $A,B\subseteq\R$ be non-empty and bounded above. Prove that
	$\sup(A+B)=\sup A+\sup B$, where $A+B=\{a+b : a\in A,\ b\in B\}$. Show by
	example that $\sup(A\cdot B)=\sup A\cdot\sup B$ can fail, and identify a
	sign condition under which it holds.
\end{exercise}
